\section{Label Readers and Holes}
\label{sec:holes}

The clash tests now become a finite reader: rejection rules out the
single-cut events, while admission records the two-cut events. We show that
the resulting readers have the rejected-path property, which is the only
subtype-specific input to the word bounds.

\subsection{Words, Cuts, and Comparisons}

We use the word and cut notation of \cref{sec:words}; put $m=|w|$. Recall
from \cref{def:comparison} that a comparison of $w$ is a pair of cuts $s<e$
with $\wsuf{w}{e}\pre\wsuf{w}{s}$; its period is $\varpi=e-s$, its root is
$\wsub{w}{s}{e}$, and it is terminal when $e=m$. Comparisons are what the
readers below test.

Equivalently, $(s,e)$ is a comparison iff the suffix $\wsuf{w}{s}$ has period~$\varpi$, \ie
$w[i] = w[i-\varpi]$ for all $e \le i < m$, iff $\wsuf{w}{s}$ is a prefix of a power of its
root (\cref{lem:periodic}). Every pair $(s,m)$ is a terminal comparison. Over $\{0,1\}$, the
nonterminal comparisons of $w = 0100$ are $(0,3)$, because $\wsuf{w}{3} = 0$ is a prefix of
$w$, which has period~$3$, and $(2,3)$, because $0 \pre 00 = \wsuf{w}{2}$; the pair $(0,2)$ is
not a comparison, because $00$ is not a prefix of $0100$. Comparisons are inherited by
prefixes: if $(s,e)$ is a comparison of $w$ and $e \le h \le m$, then $(s,e)$ is a
comparison of $\wpre{w}{h}$. Extending a word may destroy a comparison when a new letter breaks its
period, and may create comparisons that end at the new cuts.


\begin{definition}[Reader hole]\label{def:reader-hole}
A \emph{reader} $(\mathcal{M},\Adm,\Rej)$ consists of a deterministic finite automaton
$\mathcal{M} = (Q,q_0,\delta)$ over a finite alphabet $\Gamma$, with transition function
$\delta : Q \times \Gamma \to Q$ extended to words by $\delta(q,\eps) = q$ and
$\delta(q,wa) = \delta(\delta(q,w),a)$, together with a relation $\Adm \subseteq Q \times Q$,
the \emph{admission relation}, and a set $\Rej \subseteq Q$ of \emph{rejected} states. A
state is \emph{reachable} if it is $\delta(q_0,w)$ for some word $w$. For a word $w$, write
$q_i = \delta(q_0,\wpre{w}{i})$ for the state at cut $i$. A word $w$ is a \emph{hole} of
$(\mathcal{M},\Adm,\Rej)$ if $q_{\len{w}} \in \Rej$ and no comparison $(s,e)$ of $w$ satisfies
$\Adm(q_s,q_e)$. States outside $\Rej$ are accepted.
\end{definition}
Rejection identifies a candidate counterexample word. An admitted comparison
disqualifies that candidate. Thus a hole requires both a rejected final state and the
absence of every admitted comparison.



\subsection{The Label Reader}\label{sec:label-reader}
For a set of signed variables define its one-letter updates

\begin{equation}\label{eq:g37}
\begin{aligned}
 E_i(U)&=\{v:\exists u\in U\quad\Up(i,u,v)\},\\
 G_i(L)&=\{v:\exists u\in L\quad\Lo(i,v,u)\}.
 \end{aligned}
\end{equation}

For a flip-closed, sign-coherent system $\psi$ and positive query roots
$x^+,y^+$, the \emph{label reader} has states $(U,L,b)$, where $U,L\subseteq V^{\pm}$
and $b\in\{0,1\}$. Its initial state and transition are

\begin{equation}\label{eq:g38}
q_0=(D_\psi\mathord\uparrow x^+,D_\psi\mathord\downarrow y^+,0),
 \qquad
 \delta_i(U,L,b)=(E_iU,G_iL,b\lor\ready(U,L)).
\end{equation}

The bit $b$ is a \emph{latch}: once set, it stays set and records readiness
at an earlier cut. It is needed because readiness at any earlier cut
excludes every extension from being unsafe, even if readiness later disappears.
Since $|V^\pm|=2k$, each of $U,L$ has $2^{2k}$ possible values;
together with the two latch values this gives
$2\cdot2^{2k}\cdot2^{2k}=2\cdot4^{2k}$ states.
All transitions and tests are computable
from the finite closure. Write $q_j$ for the state after $w[:j]$.
The path rules (\cref{lem:path-rules}) and induction on $j$ prove the run invariant

\begin{equation}\label{eq:g39}
q_j=(U_j,L_j,b_j),\qquad
 b_j=1\iff\exists h<j\quad\ready(U_h,L_h).
\end{equation}

Designate as rejected the states

\begin{equation}\label{eq:g40}
\Rej=\{(U,L,b):b=0\land\neg\ready(U,L)
                         \land\neg\childtest(U)\},
\end{equation}

and define the binary state relation

\begin{equation}\label{eq:g41}
\Adm((U,L,b),(U',L',b'))
 \iff L\cap U'\ne\varnothing\ \lor\ U\cap\flip U'\ne\varnothing.
\end{equation}

A word $w$ is a hole precisely when
\begin{equation}\label{eq:g42}
q_m\in \Rej\quad\land\quad
 \forall s<e\le m,\quad w[e:]\preceq w[s:]\Longrightarrow\neg \Adm(q_s,q_e).
\end{equation}
The set of holes of a reader need not be a regular language, even for a
label reader, since the comparison test relates suffixes that start at
different cuts; \cref{sec:general-holes} gives an example. The argument
does not require regularity.

There is one reader construction, with two instances: the \emph{left and
right label readers}. The left label
reader uses $\psi$ with the roots $(x^+,y^+)$; the right one uses the same
system $\psi$ with the roots $(y^-,x^-)$, which are positive for the
complementary sign map (\cref{sec:spine}). Closure, path updates, readiness,
child, and admission are the same for both; only the initial state differs.
By \cref{lem:fixed-witness}, a top prefix of $y^-$ is a bottom prefix of
$y$, so the right label reader tests right unsafety.
Write $(\mathcal M_\theta,\Adm_\theta,\Rej_\theta)$ for the instance on
$\psi$ with the roots $(X_\theta,Y_\theta)$ defined in
\cref{sec:fixed-witness}.

\begin{example}[Reading the running query]\label{ex:running-reader}
We run the left-side construction on the signed system of
\cref{ex:running-signed} along $00$, for the failing query $x\le y$ and,
for contrast, for the entailed query $y\le x$, which exchanges the roots.
Neither half is the right label reader of $x\le y$, which would start
from the roots $(y^-,x^-)$. The last three columns give the latch, readiness,
and acceptance status at each cut, before the next letter is read.
\begin{center}
\small
\begin{tabular}{@{}ccccccl@{}}
\toprule
Query & cut $j$ & $U_j$ & $L_j$ & $b_j$ & $\ready$ & status\\
\midrule
$x\le y$ & $0$ & $\{x^+\}$ & $\{y^+\}$ & $0$ & false & accepted (child)\\
 & $1$ & $\{a^-\}$ & $\varnothing$ & $0$ & false & rejected\\
 & $2$ & $\varnothing$ & $\varnothing$ & $0$ & false & rejected\\
\midrule
$y\le x$ & $0$ & $\{y^+\}$ & $\{x^+\}$ & $0$ & false & accepted (child)\\
 & $1$ & $\{a^-\}$ & $\{a^-\}$ & $0$ & true & accepted (ready)\\
 & $2$ & $\varnothing$ & $\varnothing$ & $1$ & false & accepted (latch)\\
\bottomrule
\end{tabular}
\end{center}
For $x\le y$, the word $0$ is a hole: its final state is rejected,
and its only comparison $(0,1)$ is not admitted. The word $00$ is also
a hole. The empty word is accepted by the child test; it is not a
left-unsafe word. For the entailed reverse query, readiness at cut $1$
sets the latch at cut $2$, so the empty upper and lower sets there do
not cause rejection.

Rejection alone does not certify a hole. In the cross event of
\cref{ex:cross-event}, the run on $11$ has $b_j=0$ throughout and ends
in $(\{b^+\},\varnothing,0)$, a rejected state. Nevertheless $(0,1)$
is cross-admitted, so $11$ is not a hole.
\end{example}

\subsection{The Rejected-Path Property}\label{sec:interface}
For a general reader $(\mathcal{M},\Adm,\Rej)$, write $q_j=\delta(q_0,w[:j])$. If $q_m\in\Rej$,
let $J=J(w)$ be the least cut with $q_J\in\Rej$, and put $r(w)=|w|-J(w)$. The reader has the
\emph{rejected-path property} if it satisfies \emph{betweenness} and
the \emph{admission cone}:

\begin{description}[style=nextline,leftmargin=1em,font=\normalfont\itshape]\item[Betweenness] if two prefixes of a word end in $\Rej$, every intermediate
  prefix also ends in $\Rej$.
\item[Admission cone] in a word ending in $\Rej$, every comparison
  $(s,e)$ that satisfies $\Adm(q_s,q_e)$ has $e\le J$.

\end{description}

\begin{lemma}[Rejected paths of label readers]\label{lem:between}\label{lem:cone}
The label reader, and hence both the left and the right label reader, satisfies the rejected-path
property, without any satisfiability assumption.
\end{lemma}

\begin{proof} Rejection at a cut implies $\neg\childtest(U)$.
Every reachable $U$ is upward closed under $D_\psi$: this holds at
$q_0$ and is preserved by each $E_i$, using \eqref{eq:g22} and transitivity.
If $E_iU\ne\varnothing$, \eqref{eq:g22} supplies an upper constructor with root
$z$ reached from some $u\in U$. Upward closure puts $z\in U$,
giving $\childtest(U)$, a contradiction. Thus $U$ becomes
empty on the next letter and remains empty thereafter. Both disjuncts of \eqref{eq:g41} require a
nonempty upper set at their second cut. Any admission therefore has
$e\le J$, proving the admission cone.

For betweenness, after a rejected cut every later upper set is empty. A later
state can cease to be rejected only through readiness or an already set
latch; a child can no longer occur. Once readiness occurs, it is present
at that cut and sets the latch at the next one, making every subsequent
state accepted. Consequently it cannot occur between two rejected cuts.
This proves betweenness for either instance of the construction. \end{proof}

The proof gives more than the admission cone: it excludes later
admissions even for pairs that are not comparisons. The rejected-path property retains
only what the tail proof needs; this stronger restriction may help sharpen
subtype-specific bounds. Rejection need not be absorbing: a later ready
state can become accepted permanently.

\subsection{Unsafe Words and Holes}\label{sec:side-bridge}
\begin{theorem}[Unsafe words are holes]\label{thm:holes}
If $\Sat_c(\phi)$, then for either side $\theta$,

\begin{equation}\label{eq:g44}
w\text{ is }\theta\text{-unsafe}
 \quad\Longleftrightarrow\quad
 w\text{ is a hole of }(\mathcal{M}_\theta,\Adm_\theta,\Rej_\theta).
\end{equation}\end{theorem}

\begin{proof} By \cref{thm:four-spine}, unsafety is equivalent to absence
of a label clash in the side's four-spine extension. Satisfiability excludes
inherited clashes in $\psi$, so \cref{thm:clashes} identifies
this with the simultaneous negation of \eqref{eq:g36}.
The run and test formulas \eqref{eq:g39}--\eqref{eq:g41} identify those negations precisely
with \eqref{eq:g42}, using the right label reader for the right side. Every implication is reversible. The argument uses
no positive-arity assumption; the letter-free child test also handles the
empty word over the empty alphabet. \end{proof}

\subsection{Consequences for General Readers}
For any reader satisfying the rejected-path property, the following consequences hold.

\begin{corollary}[Hole criterion]\label{cor:normal}
A rejected word $w$ is a hole iff it has no admitted comparison $(s,e)$ with $e \le J(w)$.
\end{corollary}

\begin{lemma}[Rejected extension]\label{lem:extension}
If $v$ is a hole and $w$ is a rejected word with $v \pre w$, then $w$ is a hole.
\end{lemma}

\begin{proof}
Since $v$ is rejected, $J(w) = J(v) \le \len{v}$. An admitted comparison $(s,e)$ of $w$ has
$e \le J(w) \le \len{v}$ by the admission cone, so it is also a comparison of $v$, with the
same cut states. This contradicts the assumption that $v$ is a hole.
\end{proof}

\Cref{tab:reader-hypotheses} separates the subtype-specific rejected-path property
from the general word arguments.
\begin{table}[tbp]
\caption{Hypotheses used by the word bounds.}
\label{tab:reader-hypotheses}
\centering
\begin{tabular}{@{}p{.44\linewidth}p{.51\linewidth}@{}}
\toprule
Result & Hypotheses on the reader and word\\
\midrule
Rejected extension (\cref{lem:extension}) & Admission cone; a hole and a rejected extension\\
Bounded tail (\cref{thm:brt}) & Admission cone and betweenness; some hole\\
Return gap and depth (\cref{lem:return-gap,thm:depth}) & Any reader; a rejected word, nonempty for depth\\
Bounded-depth compression (\cref{thm:search}) & Any reader; a nonempty hole of bounded depth\\
\bottomrule
\end{tabular}
\end{table}

\begin{example}[Loop deletion can create an admission]\label{ex:loop-deletion}
Consider the reader on $\{0,1\}$ with initial state $q$, rejected set
$\Rej=\{r\}$, sole admitted pair $(q,b)$, and all transitions as follows:
\begin{center}
\begin{tabular}{@{}c|cccccc@{}}
state & $q$ & $a$ & $b$ & $c$ & $r$ & $h$\\
\hline
$0$ & $a$ & $a$ & $c$ & $h$ & $r$ & $h$\\
$1$ & $h$ & $b$ & $h$ & $r$ & $r$ & $h$
\end{tabular}
\end{center}
Here $h$ is an accepted sink. State $q$ occurs only at cut $0$, so a
comparison is admitted only if it starts at cut $0$ and ends at a cut in
state $b$. Deleting the loop $a\xrightarrow{0}a$ between the equal-state
cuts $1$ and $2$ of $00101$ gives $0101$; both runs end in the rejected
state $r$:
\[
\begin{array}{r@{\quad}l}
 00101: & q\xrightarrow{0}a\xrightarrow{0}a\xrightarrow{1}b
          \xrightarrow{0}c\xrightarrow{1}r,\\[2pt]
 0101: & q\xrightarrow{0}a\xrightarrow{1}b\xrightarrow{0}c\xrightarrow{1}r.
\end{array}
\]
In $00101$, state $b$ occurs only at cut $3$, so the only candidate is
$(0,3)$. It is not a comparison: the suffix $01$ at cut $3$ is not a
prefix of $00101$. Thus $00101$ is a hole. In $0101$, state $b$ occurs at
cut $2$, and the candidate $(0,2)$ is a comparison, since the suffix $01$
at cut $2$ is a prefix of $0101$. Its states are $(q,b)$, so it is
admitted, and the shortened word is not a hole.

This reader has the rejected-path property. State $r$ is absorbing,
which proves betweenness. Any admission must start at $q$ and end at
$b$; $q$ occurs only initially, and $b$ can occur only before the first
visit to $r$. Hence every admitted comparison of a rejected word ends
before $J$, proving the admission cone. This example motivates the
simultaneous comparison summaries in \cref{sec:search}.
\end{example}
